
\documentclass{article}
\usepackage{colortbl}
\usepackage{makecell}
\usepackage{multirow}
\usepackage{supertabular}

\begin{document}

\newcounter{utterance}

\centering \large Interaction Transcript for game `cladder', experiment `full\_v1.5\_default', episode 9064 with Qwen3.5{-}27B.
\vspace{24pt}

{ \footnotesize  \setcounter{utterance}{1}
\setlength{\tabcolsep}{0pt}
\begin{supertabular}{c@{$\;$}|p{.15\linewidth}@{}p{.15\linewidth}p{.15\linewidth}p{.15\linewidth}p{.15\linewidth}p{.15\linewidth}}
    \# & \multicolumn{2}{c}{Player} && \multicolumn{2}{c}{Game Master} \\
    \hline

    \theutterance \stepcounter{utterance}  
    & & & \multicolumn{4}{p{0.6\linewidth}}{
        \cellcolor[rgb]{0.9,0.9,0.9}{
            \makecell[{{p{\linewidth}}}]{
                \texttt{\tiny{[P1$\langle$GM]}}
                \texttt{You are an expert in causal inference. The following question is not a typical commonsense query, but rather a meticulously designed question created by a professor specializing in causal inference, intended to assess the students' mastery of the course content.} \\
\texttt{Imagine a self{-}contained, hypothetical world with only the following conditions, and without any unmentioned factors or causal relationships: Pexu has a direct effect on hwax. Hwax has a direct effect on rukz. Method 1: We look at how pexu correlates with rukz case by case according to hwax. Method 2: We look directly at how pexu correlates with rukz in general. To understand how pexu affects rukz, is it more correct to use the Method 1 than Method 2?} \\
\texttt{Start your answer with 'Yes' or 'No', followed by additional reasoning or evidence to support your explanation.} \\
            }
        }
    }
     \\ \\

    \theutterance \stepcounter{utterance}  
    & \multicolumn{4}{p{0.6\linewidth}}{
        \cellcolor[rgb]{0.561,0.737,0.561}{
            \makecell[{{p{\linewidth}}}]{
                \texttt{\tiny{[P1$\rangle$GM]}}
                \texttt{No} \\
\\ 
\texttt{In the described hypothetical world, the causal structure is a simple chain: \$Pexu \textbackslash{}rightarrow Hwax \textbackslash{}rightarrow Rukz\$. Here, \$Hwax\$ acts as a **mediator** (or intermediate variable) through which \$Pexu\$ exerts its effect on \$Rukz\$.} \\
\\ 
\texttt{**Method 2** (looking at the general correlation between \$Pexu\$ and \$Rukz\$) estimates the **Total Effect** of \$Pexu\$ on \$Rukz\$. This includes both the direct path (which is zero in this specific chain) and the indirect path through \$Hwax\$. Since the question asks how \$Pexu\$ affects \$Rukz\$ in general, the total effect is the correct quantity to estimate.} \\
\\ 
\texttt{**Method 1** (stratifying or conditioning on \$Hwax\$) attempts to estimate the **Direct Effect** of \$Pexu\$ on \$Rukz\$ while holding \$Hwax\$ constant. However, because \$Hwax\$ is a mediator, conditioning on it blocks the causal pathway from \$Pexu\$ to \$Rukz\$. In a pure mediation chain where there is no direct effect, conditioning on the mediator will result in an estimated effect of zero (or near zero), effectively "adjusting away" the very causal mechanism you are trying to measure.} \\
\\ 
\texttt{Therefore, to understand the overall causal impact of \$Pexu\$ on \$Rukz\$, Method 2 is correct, while Method 1 is incorrect because it introduces **over{-}adjustment bias** (specifically, blocking the causal path).} \\
            }
        }
    }
    & & \\ \\

    \theutterance \stepcounter{utterance}  
    & & & \multicolumn{2}{p{0.3\linewidth}}{
        \cellcolor[rgb]{0.9,0.9,0.9}{
            \makecell[{{p{\linewidth}}}]{
                \texttt{\tiny{[GM$|$GM]}}
                \texttt{No} \\
            }
        }
    }
    & & \\ \\

    \theutterance \stepcounter{utterance}  
    & & & \multicolumn{2}{p{0.3\linewidth}}{
        \cellcolor[rgb]{0.9,0.9,0.9}{
            \makecell[{{p{\linewidth}}}]{
                \texttt{\tiny{[GM$|$GM]}}
                \texttt{no} \\
            }
        }
    }
    & & \\ \\

    \theutterance \stepcounter{utterance}  
    & & & \multicolumn{2}{p{0.3\linewidth}}{
        \cellcolor[rgb]{0.9,0.9,0.9}{
            \makecell[{{p{\linewidth}}}]{
                \texttt{\tiny{[GM$|$GM]}}
                \texttt{game\_result = WIN} \\
            }
        }
    }
    & & \\ \\

\end{supertabular}
}

\end{document}
